% ============ 附录 B 进一步学习的建议 (原书 pp.415-419) ============
\chapter*{附录 B: 进一步学习的建议}

代数数论的经典著作有:

[1] Gauss, C. F., \emph{Disquisitiones Arithmeticae}, 1801（英译本 Springer-Verlag, 1986）;

[2] Hilbert, D., \emph{Zahlbericht}, 1897（英译本 \emph{The Theory of Algebraic Number Fields}, Springer-Verlag, 1998）;

[3] Hecke, E., \emph{Vorlesungen \"uber die Theorie der algebraischen Zahlen}, 1923（英译本 \emph{Lectures on the Theory of Algebraic Numbers}, GTM 77, Springer-Verlag, 1981）;

[4] Weil, A., \emph{Basic Number Theory}, 3rd, Springer-Verlag, 1995（第 $1$, $2$ 版于 $1967$, $1973$ 年）.

Gauss [1] 建立了二次域的基本理论, 为代数数论的开端. Hilbert [2] 的起因是德国数学会于 $1893$ 年邀请 Hilbert 写一份关于代数数论研究现状的报告, 结果 Hilbert 写成一本研究专著, 不仅详细总结和整理了 $19$ 世纪代数数论成果, 而且在分歧理论等许多方面有系统的发展. Hecke [3] 给出 Dedekind zeta 函数的基本解析特性, 包括详细证明了它的函数方程, 书中反映了作者关于模形式（theta 函数）的思想. Weil [4] 是现代代数数论的奠基性著作.

关于代数数论的经典内容, 已有许多著作, 其中有:

[5] Borevich, Z. I., Shafarevich, I. R., \emph{Number Theory}, 3rd ed., Academic Press, 1985（原书俄文, 1964）;

[6] Weiss, E., \emph{Algebraic Number Theory}, McGraw-Hill, 1963;

[7] Marcus, D. A., \emph{Number Fields}, Springer-Verlag, 1977;

[8] Ireland, K., Rosen, M., \emph{A Classical Introduction to Modern Number Theory}, 3rd, GTM 84, Springer-Verlag, 1986;

[9] Neukirch, J., \emph{Algebraic Number Theory}, Springer-Verlag, 1999;

[10] Hasse, H., \emph{Zahlentheorie}, Academic Press, 1949（英译本 1964）;

[11] Narkiewicz, W., \emph{Elementary and Analytic Theory of Algebraic Numbers}, Warsaw, Polish Sci. Publ., 1974;

[12] Artin, E., \emph{Algebraic Numbers and Algebraic Functions}, Gordon \& Breach, 1967;

[13] Lang, S., \emph{Algebraic Number Theory}, 2nd, GTM 110, Springer-Verlag, 1994;

[14] Esmonde, J., Murty, M. R., \emph{Problems in Algebraic Number Theory}, GTM 190, Springer-Verlag, 1999.

其中 [5], [6], [7], [8] 和 [12] 为好的入门书. [5] 以不定方程为中心议题讲述代数数论. [8], [10] 和 [12] 同时讲述数域和函数域, [11] 对数域的代数和解析理论介绍全面, 并包含许多研究问题, 书末有丰富的文献目录. 对代数数论经典内容有基本了解之后建议读 [9] 和 [13]. 而 [14] 中收集了五百余个代数数论的习题和答案.

关于 $20$ 世纪近代代数数论, 建议读者先看

[15] Cassels, J. W. S., Fr\"ohlich, A., \emph{Algebraic Number Theory}, Academic Press, 1967,

这本文集由各家撰文介绍代数数论（整体域和局部域理论）, 类域论、模形式、椭圆曲线、伽罗瓦表示的基本知识和方法, 对于近代数论可以有一个基本了解. 讲述 $\Q$ 上二次型代数理论（Hasse 局部--整体原则）的有

[16] Serre, J.-P., \emph{A Course in Arithmetic}, 1973（中译本, 上海科技出版社, 1980）;

[17] Cassels, J. W. S., \emph{Rational Quadratic Forms}, Academic Press, 1978.

关于 $p$-adic 分析和局部域理论的专门著作有

[18] Koblitz, N., \emph{$p$-adic Numbers, $p$-adic Analysis and Zeta Functions}（第一版）, GTM 58, Springer-Verlag, 1984;

[19] Cassels, J. W. S., \emph{Local Fields}, Cambridge Univ. Press, 1986;

[20] Serre, J.-P., \emph{Local Fields}, GTM 67, Springer-Verlag, 1979.

其中 [18] 通俗易懂并且注重介绍思想. [20] 的内容较多. 关于类域论的专门著作有:

[21] Neukirch, J., \emph{Class Field Theory}, Springer-Verlag, 1986;

[22] Hasse, H., \emph{Klassenk\"orper}, Marburg, 1933;

[23] Artin, E., Tate, J., \emph{Class Field Theory}, Benjamin, 1968;

[24] Iwasawa, K., \emph{Local Class Field Theory}, Oxford Univ. Press, 1986.

其中 [21] 是学习类域论最好的入门书, [22] 和 [23] 是类域论经典著作, [24] 是局部类域论的好书. 此外, Weil [4] 用现代观点讲述类域论.

关于分圆域近代理论有两本专门著作:

[25] Washington, L. C., \emph{Introduction to Cyclotomic Fields}, 2nd, GTM 83, Springer-Verlag, 1996;

[26] Lang, S., \emph{Cyclotomic Fields}, 2nd ed., GTM 121, Springer-Verlag, 1990.

关于模形式的专门著作有

[27] Miyake, T., \emph{Modular Forms}, Springer-Verlag, 1989;

[28] Lang, S., \emph{Introduction to Modular Forms}, Springer-Verlag, 1976;

[29] Ogg, A. P., \emph{Modular Forms and Dirichlet Series}, Benjamin, 1965;

[30] Weil, A., \emph{Dirichlet Series and Automorphic Forms}, Lecture Notes in Math. 189, Springer-Verlag, 1971.

关于椭圆曲线算术理论目前也有许多著作. 由于它与模形式理论联系密切, 这些书中也包括模形式的相关内容:

[31] Silverman, J. H., \emph{The Arithmetic of Elliptic Curves}, GTM 106, Springer-Verlag, 1986;

[32] Koblitz, N., \emph{Introduction to Elliptic Curves and Modular Forms}, 2nd, GTM 97, Springer-Verlag, 1984;

[33] Husem\"oller, D., \emph{Elliptic Curves}, GTM 111, Springer-Verlag, 1987;

[34] Knapp, A. W., \emph{Elliptic Curves}, Princeton Univ. Press, 1992;

[35] Silverman, J. H., \emph{Advanced Topics in the Arithmetic of Elliptic Curves}, GTM 151, Springer-Verlag, 1994.

其中 [31] 是学习椭圆曲线算术理论的理想入门书. 由于这一理论近年来进步很快, Silverman 又写了后续著作 [35]. 而 [32] 是以一个著名数学问题（Congruent numbers）为线索介绍椭圆曲线和模形式理论的一些重要方面的.

现代算术几何和自守表示理论需要代数几何、微分几何、李群和李代数、群的无限维表示和调和分析等多方面知识, 这方面的著作有:

[36] Cornell, G., Silverman, J. H., \emph{Arithmetic Geometry}, Springer-Verlag, 1986;

[37] Borel, A., Casselman, W., \emph{Automorphic Forms, Representations and $L$-functions}, AMS Proceedings of Symposia in Pure Mathematics 33, 1979;

[38] Coates, J., Yau, S.-T., \emph{Fermat's Last Theorem}, International Press, 1997;

[39] Cornell, G., Silverman, J. H., Stevens, G., \emph{Modular Forms and Fermat's Last Theorem}, Springer-Verlag, 1997;

[40] Iwaniec, H., \emph{Introduction to the Spectral Theory of Automorphic Forms}, Revista Matem\'atica Iberoamericana, Madrid, 1995;

[41] Gelbart, S., \emph{Automorphic Forms on Adele Groups}, Princeton Univ. Press, 1975;

[42] Shimura, G., \emph{Introduction to the Arithmetic Theory of Automorphic Functions}, Princeton Univ. Press, 1971;

[43] Terras, A., \emph{Harmonic Analysis on Symmetric Spaces and Applications}, I and II, Springer-Verlag, 1985 and 1987;

[44] Jacquet, H., Langlands, R. P., \emph{Automorphic Forms on GL(2)}, Lecture Notes in Math. 114, Springer-Verlag, 1970;

[45] Ramakrishnan, D., Valenza, R. J., \emph{Fourier Analysis on Number Fields}, Springer-Verlag, 1998;

[46] Goss, D., \emph{Basic Structures of Function Field Arithmetic}, Springer-Verlag, 1996.

其中 [36] 和 [37] 分别是关于算术几何与自守表示的文集, 可用来对这两个领域作基本的了解. $1993$ 年和 $1995$ 年, 在香港中文大学和美国 Boston 大学举行了会议讨论和学习 A. Wiles 关于费马猜想的证明, [38] 和 [39] 是这两次会议的文集, 从中可看到证明中需要的现代数论和算术几何工具. [40] 和 [45] 是最近出版的学习自守表示的入门书, 其中 [40] 偏重于解析方面, [45] 讲述现代数论所需要的调和分析工具. Drinfeld 在证明函数域上二维局部 Langlands 猜想中建立了 Drinfeld 模理论, [46] 总结了近年来在这方面的研究成果.

最后介绍几本数论及其应用方面的著作:

[47] Koblitz, N., \emph{A Course in Number Theory and Cryptography}, 2nd, GTM 114, Springer-Verlag, 1994;

[48] Moreno, C., \emph{Algebraic Curves over Finite Fields}, Cambridge Univ. Press, 1991;

[49] Stichtenoth, H., \emph{Algebraic Function Fields and Codes}, Springer-Verlag, 1993;

[50] Li, W.-C. W., \emph{Number Theory with Applications}, World Scientific, 1996;

[51] Cohen, H., \emph{A Course in Computational Algebraic Number Theory}, GTM 138, Springer-Verlag, 1996.

其中 [47] 是数论在密码学中的应用, [48], [49] 讲述有限域上代数曲线算术理论及在纠错码方面的应用（代数几何码）, [50] 讲述函数域和自守形式理论及在特征和估计和图论（Ramanujan 图）的应用, 而 [51] 是计算代数数论的最近著作.
